\begin{pastexam}[label={px:gk3-100-ds-07}]{高考三級資訊處理 100 年 第七題}
（申論；10 分）

將數字 25,5,75,0,60,10,55,15,45,15 依序存入一維陣列如下，以 heap sort 方式進行由小到大的排序。請顯示其在
第一次執行完 initial heap 步驟後的一維陣列內容。（10 分）
\begin{center}
\begin{tabular}{|l|*{10}{l|}}
\hline
index & 0 & 1 & 2 & 3 & 4 & 5 & 6 & 7 & 8 & 9 \\ \hline
data & 25 & 5 & 75 & 0 & 60 & 10 & 55 & 15 & 45 & 15 \\ \hline
\end{tabular}
\end{center}
\end{pastexam}

\begin{pastsoln}
答案：
\begin{center}
\begin{tabular}{|l|*{10}{c|}}
\hline
index & 0 & 1 & 2 & 3 & 4 & 5 & 6 & 7 & 8 & 9 \\ \hline
data & 75 & 60 & 55 & 45 & 15 & 10 & 25 & 15 & 0 & 5 \\ \hline
\end{tabular}
\end{center}

由小到大排序要建\textbf{最大堆積}：每次把堆積頂端（最大值）換到陣列尾端。 陣列從 $0$ 編號，節點 $i$ 的
子節點是 $2i+1$、$2i+2$。 建初始堆積用由下而上的方法（資結書第~8.2~節、虛擬碼~8.3）：從最後一個內部節點
$\lfloor n/2 \rfloor - 1 = 4$ 往回到 $0$，對每個節點做向下調整——和較大的子節點比，比它小就交換，
一路往下直到不必換。

\begin{itemize}
\item $i = 4$（60）：子節點 $9$ 是 15，$60 \ge 15$，不動。\\
陣列：25, 5, 75, 0, 60, 10, 55, 15, 45, 15
\item $i = 3$（0）：子節點 15、45，和較大的 45 交換。\\
陣列：25, 5, 75, 45, 60, 10, 55, 15, 0, 15
\item $i = 2$（75）：子節點 10、55，$75$ 最大，不動。\\
陣列：25, 5, 75, 45, 60, 10, 55, 15, 0, 15
\item $i = 1$（5）：子節點 45、60，和 60 交換；5 到了 $4$，子節點 15，再交換。\\
陣列：25, 60, 75, 45, 15, 10, 55, 15, 0, 5
\item $i = 0$（25）：子節點 60、75，和 75 交換；25 到了 $2$，子節點 10、55，再和 55 交換。\\
陣列：75, 60, 55, 45, 15, 10, 25, 15, 0, 5
\end{itemize}
檢查：每個節點都不小於它的子節點（$75 \ge 60, 55$；$60 \ge 45, 15$；$55 \ge 10, 25$；$45 \ge 15, 0$；$15 \ge 5$），
是最大堆積。
\end{pastsoln}
